\section{Measuring: Compiler-Introduced and Compiler-Removed Channels}
\label{sec:measure}

Given one kernel, one lowering pipeline and one optimisation level, is the binary's behaviour distinguishable across two classes of secret input, and if so, through which operand did the secret get there. The answer is a measurement, not a proof: it holds for the two classes it exercised, on the toolchain that built the binary, and it is reported with those coordinates attached. \S\ref{sec:evaluation:select} and \S\ref{sec:evaluation:structure} interpret the mechanisms; this section states the method and the full results.

\subsection{Methodology}

The rig builds one binary per cell, where a cell is one kernel lowered through one MLIR pipeline at one LLVM optimisation level, and runs it on two classes of secret input under two Valgrind tools. Callgrind, started and stopped around the kernel call, counts the instructions the simulated processor executed (Ir), the conditional branches it executed (Bc) and its write references (Dw). Memcheck, with the secret buffer marked undefined before the call, tracks definedness through every arithmetic and memory operation and reports where an undefined value reaches a conditional jump or move (the control channel, cf) or the address of a load or store (the address channel, addr). This is Langley's ctgrind construction~\cite{langley2010ctgrind}: Memcheck's definedness shadow bits serve as taint tags. Counts say that two runs did different amounts of work; taint says which operand carried the secret to an observation, and needs no magnitude.

Taint is primary and counts corroborate. A cell is distinguishable if Memcheck reports a channel or if a count differs between the classes, and a count difference on its own is admitted only past two guards, because the count channel has a confound that is not noise. Valgrind disables address-space randomisation, so memory layout is a deterministic function of a run's arguments and environment, and byte-identical secrets stored at file paths of different length execute counts that differ by several hundred instructions; repeating at one fixed path shows a spread of zero and hides this. Each class is therefore measured at five contexts of varying path length, paired by context, and a count difference is reported only if its median exceeds the largest spread either class shows across contexts (the floor) and the paired differences agree across contexts (stability). Layout manufactures false positives, never false zeros, so a count that fails a guard is uninterpretable rather than clean, and the layout-insensitive taint channel decides. Two consequences follow. Memcheck's control report covers a conditional move as well as a jump; a comparison alone does not necessarily trigger that diagnostic. Where this distinction matters, the verdict is taken on the address channel. And a leak whose only signature is the address touched is invisible to the counts and visible to Memcheck alone.

Secret classes are chosen, not sampled: each pair sits on opposite sides of the mechanism its kernel probes. Comparing two cells gives four outcomes. Relative to a baseline cell, a second cell is oblivious if neither leaks, leak-present if both do, introduced if only the second does and removed if only the baseline does; the baseline is a parameter of the run and is named in every report, and below it is the plain lowering at -O0 unless stated.

Every number below and in Section~\ref{sec:evaluation} comes from one run on one machine: an AMD Ryzen AI 9 HX 370 under Ubuntu 24.04 with mlir-opt 18.1.3, clang 18.1.3 with \texttt{-mavx2 -mno-avx512f} (Valgrind 3.22 does not decode AVX-512), Valgrind 3.22.0, glibc 2.39, kernels of 4096 elements and 256×256 matrices. The kernels that allocate pass through glibc's CPUID-dispatched routines, so their counts may differ on another processor.

\subsection{Kernels and the compiler axis}
Five kernels probe the obligations of \S\ref{sec:verify}, each with the two classes that bracket its mechanism: \texttt{matvec}, dense with fixed bounds, the negative control (zero weights against Gaussian); \texttt{cond\_reduce}, a branch on the secret sum (all zero against all one); \texttt{mask\_select}, a branchless \texttt{arith.select} on a secret mask (all false against all true); \texttt{idx\_gather}, a load at a secret index (index 0 against random); and \texttt{dynshape}, a secret extent used as allocation size and loop bound (1 against 4096). Tensor-typed variants \texttt{matvec\_t}, \texttt{mask\_select\_t} and \texttt{dynshape\_t} give one-shot bufferization real decisions to make; \texttt{scatter} under the sparsifier, against a hand-written dense reference, is described in Section~\ref{sec:evaluation:structure}.

The MLIR axis is six pipelines that differ only before a common tail into the LLVM dialect: the plain \texttt{linalg}→\texttt{scf}→\texttt{cf} path (P0), the affine path (P1), canonicalisation and CSE ahead of P0 (P2), affine super-vectorisation (P3), one-shot bufferization for the tensor kernels (P4) and generalisation of named ops (P5). The LLVM axis is -O0, -O2 and -O3. Section~\ref{sec:evaluation:boundaries} sets the IR-analysis verdict beside these P0 measurements; the analysis was not swept over the other pipelines.


\subsection{Results}

Table~\ref{tab:measured-cells} gives every cell of that run. The mechanisms behind its rows are the subject of \S\ref{sec:evaluation:select} (the select) and \S\ref{sec:evaluation:structure} (shape, bufferization, sparsity); two regularities of the table as a whole are stated here.


The MLIR axis moves magnitude; the LLVM axis moves verdicts. Across the five pipelines at -O0 no kernel's verdict or channel changes. Canonicalisation ahead of lowering (P2) makes \texttt{dynshape} execute 12,447 fewer instructions and perform 4,113 fewer write references in the class-B-minus-class-A difference, and reports the same channel with the same taint. Every verdict that changes in the table changes along the -O axis, and every channel that survives -O3 is one the program computes from its input: a runtime trip count or an address.

The count channel is blind to address leaks. \texttt{idx\_gather} and the sparsified \texttt{scatter} have identical counts in both classes at every cell and are reported by Memcheck alone. A rig that measured only work (instructions, branches, cycles) would return oblivious for both.

\subsection{A production compiler}\label{sec:production}

Before the MLIR corpus we ran the same methodology on PyTorch's Inductor backend (torch 2.13, CPU), eager against compiled, on an activation-function corpus of 512×512 inputs. In exp, with classes of 0.5 against 100.0 bracketing the overflow handling of a scalar exponential, the eager build executes 35,061,760 more instructions and 3,416,064 more branches for the large class and Memcheck traces the secret to a branch condition; the compiled build shows no difference and no report. The verdict is compiler-removed. The corpus was designed to probe a scalar libm expf in the eager kernel against a branchless vectorised polynomial in Inductor's generated code. \texttt{relu} is oblivious in both builds; \texttt{gelu}'s sweep did not complete; softmax's probe was degenerate (uniform rows) and is retained as a retired result.

Freezing with constant folding writes a scale derived from \texttt{max|w|} into the generated code as a literal. Running Memcheck on the compilation itself with the weight bytes marked undefined gives 184 reports at the float-formatting site of the code generator in the frozen configuration and none in the unfrozen control, and the emitted literal equals \texttt{max|w|/127} to float32 precision across seven perturbations of the weights. Nothing was executed and nothing was timed; this is the leak the threat model files under persistence. Autotuning is a clean negative: max-autotune selects the same GEMM kernel and emits generated code identical after normalisation across dense, sparse, low-rank, subnormal, small- and large-magnitude weight classes and across independent draws within a class, consistent with the tuner benchmarking on inputs of the right shape rather than on the weights.